% =====================================================================
\chapter{Le métier de l'assurance et les invariants du produit}
\label{ch:metier}
% =====================================================================

Un projet technique appliqué à un métier se juge d'abord sur sa compréhension
de ce métier. Ce chapitre expose ce que j'ai appris du secteur, comment cela a
déterminé les choix de conception, et les règles qui en découlent — règles que
le projet s'interdit de violer.

\section{Deux acheteurs, un moteur}

Les échanges menés au démarrage du stage ont fait apparaître que \og
l'assureur \fg{} n'est pas un utilisateur unique. Deux profils aux besoins
différents cohabitent.

\subsection{Le souscripteur (\emph{underwriter})}

C'est l'utilisateur principal de la version~1. Il traite des dossiers : une
adresse arrive, il doit décider \emph{accepter}, \emph{mettre à l'étude} ou
\emph{refuser}. Ses contraintes :

\begin{itemize}[leftmargin=1.4em]
  \item \textbf{rapidité} — la décision s'insère dans un flux de travail, pas
        dans une analyse hebdomadaire ;
  \item \textbf{jugement spatial} — il pense en distance, en périmètre, en
        relief. Il regarde une carte, il ne lit pas un tableau de bord ;
  \item \textbf{traçabilité} — sa décision peut être auditée. Il doit pouvoir
        dire d'où vient chaque élément.
\end{itemize}

C'est pour ce profil que la carte reste une fonctionnalité de premier plan,
et non une décoration : elle ne montre pas les couches sources, elle montre
\emph{la jointure} — l'empreinte du bâtiment colorée par le résultat du test
géométrique, accompagnée d'une fiche de vulnérabilité.

\subsection{L'actuaire (\emph{pricing})}

Profil secondaire, hors périmètre de la version~1 mais structurant. Il
construit des modèles de tarification et de provisionnement. Ses besoins sont
presque opposés :

\begin{itemize}[leftmargin=1.4em]
  \item \textbf{pas d'interface} — il veut de la donnée, pas un écran ;
  \item \textbf{volume} — il traite des portefeuilles entiers, pas une adresse ;
  \item \textbf{variables brutes} — il veut appliquer \emph{ses} formules.
\end{itemize}

Le point décisif : l'actuaire interrogé ne voulait \emph{pas} de notre score.
Il voulait les variables par aléa, par horizon et par scénario climatique,
étiquetées et propres, pour exécuter ses propres calculs. Cette réponse a
directement conduit à l'invariant central du produit.

\section{La non-stationnarité : pourquoi la donnée historique ne suffit pas}

Le raisonnement qui fonde l'intérêt du produit mérite d'être explicité, car il
n'est pas immédiat.

Un assureur possède un actif que nous n'aurons jamais : l'historique de ses
sinistres. Pour un bien déjà couvert, cet historique est un meilleur prédicteur
que n'importe quelle reconstitution à partir de sources publiques. Il serait
donc absurde de prétendre \og mieux connaître le risque actuel \fg.

Mais le risque climatique est \textbf{non stationnaire} : le futur ne ressemble
pas au passé. Un historique de sinistres, aussi volumineux soit-il, décrit des
pertes réalisées sous le climat d'hier. Il ne peut structurellement pas
indiquer comment l'exposition d'un bâtiment donné évolue sur les 25 à 75
prochaines années. Cette information doit venir de la science du climat et des
modèles de projection, pas d'un registre de sinistres.

\begin{encadre}
\textbf{L'argument de vente n'est donc pas \og plus de données \fg.} C'est \og
la partie de l'information qui ne peut pas venir de votre historique \fg.
Aucune quantité de données propriétaires ne comble ce trou précis.
\end{encadre}

C'est pour cette raison que la trajectoire temporelle — aujourd'hui, 2050,
2100 — est le cœur du produit, et non un ornement. Et c'est pour cette raison
que la donnée climatique, un temps écartée du périmètre, a dû y revenir sous
une forme honnête (voir chapitre~\ref{ch:donnees}).

\section{Les invariants du produit}

Le projet a produit un document de référence, la \emph{constitution}
(\fichier{constitution.md}), qui fixe les règles non négociables. Chaque règle
porte un mécanisme d'application : un test automatisé, une porte CI, ou une
revue humaine. En voici les principales, avec leur justification.

\subsection{Faits et provenance uniquement — pas de score}

\begin{invariant}
Aucune réponse de l'API ne peut contenir de score composite, de jugement de
risque, de recommandation, ni d'évaluation générée par IA. La graduation
fournie par la source est transmise verbatim sous forme de texte ; toute
conversion en bandes est de la coloration de présentation et n'est
\emph{jamais} sérialisée dans le contrat canonique.
\end{invariant}

La justification est double. Juridiquement, un score tiers de risque
climatique, produit par un acteur qui n'est ni l'assureur ni le régulateur,
crée une responsabilité : s'il est erroné et qu'il a servi à refuser un
contrat, qui est responsable ? Techniquement, un score se fait
\emph{reconstruire} ou \emph{refuser} : l'actuaire qui reçoit un score
non transparent finit par ignorer la boîte noire ou par la réimplémenter.

La conséquence architecturale est forte : puisque le score disparaît, le champ
\texttt{resolution} devient \emph{le seul signal de qualité} du contrat. Chaque
objet aléa doit dire à quel niveau il a été établi.

\subsection{Un niveau de résolution sur chaque fait}

\begin{invariant}
Chaque objet aléa du contrat porte \texttt{source}, \texttt{resolution} et les
champs d'attribution exigés par la Licence Ouverte~2.0. Un objet aléa sans ces
champs échoue la validation du schéma.
\end{invariant}

Trois valeurs de résolution existent, et elles ne sont pas interchangeables :

\begin{table}[H]
\centering
\small
\caption{Les trois niveaux de résolution et leur signification exacte.}
\label{tab:resolution}
\begin{tabularx}{\textwidth}{@{}l l X@{}}
\toprule
\textbf{Valeur} & \textbf{Origine} & \textbf{Signification} \\
\midrule
\code{per-building} & Test géométrique WFS &
  Le point a été testé contre un polygone source. L'affirmation porte
  réellement sur ce bâtiment. \\
\addlinespace
\code{commune-level} & Décret ou arrêté &
  Le fait est communal \emph{par nature} (zonage sismique, classement radon).
  Ce n'est pas une dégradation : c'est la résolution maximale qui existe. \\
\addlinespace
\code{commune-level-estimate} & Repli, ou RGA &
  Nous n'avons pas pu établir le fait au niveau du bâtiment. Le fait est une
  estimation communale. À traiter avec prudence. \\
\bottomrule
\end{tabularx}
\end{table}

La distinction entre les deux derniers cas est essentielle. Un aléa sismique
\code{commune-level} est \emph{correct} : le zonage est décrétal, il n'existe
pas de version par bâtiment. Un aléa inondation \code{commune-level-estimate}
signale au contraire que la source n'a pas répondu et que l'information est
dégradée. Confondre les deux ferait passer une panne pour un fait.

\subsection{Passage verbatim de la BDNB}

\begin{invariant}
Le contrat transmet les attributs BDNB sans suppression ni renommage
silencieux. Un champ indisponible est explicitement \texttt{null}, avec
l'échec enregistré dans la liste d'erreurs — jamais omis.
\end{invariant}

C'est une conséquence directe du profil actuaire : s'il doit filtrer lui-même,
il lui faut l'intégralité du champ. La curation — quels attributs montrer dans
le panneau latéral — se fait uniquement à la couche de présentation.

\subsection{Une seule forme canonique}

\begin{invariant}
Il existe exactement un schéma de réponse par ressource. Deux formes vivantes
pour la même ressource constituent un défaut.
\end{invariant}

Cette règle a une histoire : au cours du projet, deux pipelines concurrents
produisaient deux JSON différents pour la même adresse. L'un était utilisé par
l'interface, l'autre par le lot. Les divergences silencieuses entre les deux
sont exactement le genre de bug qui détruit la confiance dans un contrat de
données. La règle a été formalisée en question ouverte (OQ-1) puis tranchée.

\subsection{Le lot est N fois l'unitaire}

\begin{invariant}
\api{POST /diagnostic/batch} traite chaque adresse à travers le pipeline
unitaire identique. Aucune seconde implémentation en masse, aucun chemin de
code d'extraction versionnée.
\end{invariant}

L'intérêt est la garantie de cohérence : il ne peut pas exister de différence
entre le résultat d'un appel unitaire et celui du même appel au sein d'un lot,
hors horodatage. Cette équivalence est testée.

\subsection{Honnêteté sur le RGA}

\begin{invariant}
Tant qu'aucune source vectorielle n'est câblée, le retrait-gonflement des
argiles (RGA) est livré avec \code{resolution: "commune-level-estimate"} et
rien ne peut revendiquer une résolution par bâtiment pour cet aléa.
\end{invariant}

Le RGA est un motif majeur de sinistres en France. La tentation d'afficher un
résultat par bâtiment est donc forte. Mais la donnée disponible au moment du
stage était communale. Le choix a été de l'annoncer plutôt que de la surinterpréter.

\subsection{Pas de code dormant}

\begin{invariant}
Les fonctionnalités retirées sont supprimées, non commentées ou laissées non
référencées. L'historique Git est l'archive.
\end{invariant}

Cette règle est inconfortable mais saine : elle interdit de conserver du code
\og au cas où \fg, qui est la principale source d'ambiguïté dans un dépôt. Elle
a conduit à supprimer des milliers de lignes — connecteurs, schémas, composants
d'interface — et à assumer que Git garderait la trace.

\section{Synthèse}

\begin{table}[H]
\centering
\small
\caption{Décisions métier verrouillées lors de la phase de cadrage, et leur
traduction en règle technique.}
\label{tab:decisions-metier}
\begin{tabularx}{\textwidth}{@{}c X X@{}}
\toprule
\textbf{\#} & \textbf{Décision} & \textbf{Conséquence technique} \\
\midrule
1 & L'acheteur est l'assureur français, pilote au poste de souscription &
  Interface orientée dossier unique, pas portefeuille \\
2 & Le livrable principal est un contrat JSON, pas une interface &
  Schéma Pydantic versionné comme artefact central \\
3 & La jointure WFS $\times$ BDNB est le périmètre ; Copernicus et
  Open-Meteo hors v1 &
  Suppression des connecteurs climatiques de la version 1 \\
4 & La transaction est l'adresse unique &
  \api{POST /diagnostic/adresse} comme primitive, le lot en dérivé \\
5 & 13 aléas, certification sur 4 d'entre eux &
  Carte de couches en configuration, pas en logique \\
6 & RGA étiqueté communal &
  Résolution forcée au niveau du schéma \\
7 & Sismique et radon communaux par nature &
  Résolution permanente \code{commune-level} \\
8 & BDNB verbatim &
  Test de conformité clé par clé \\
9 & Score composite mort &
  Schéma qui refuse la sérialisation de champs de score \\
10 & Bandes D03 en présentation seulement &
  Coloration côté frontend, hors contrat \\
11 & Le champ de résolution est le signal de qualité &
  Tests de dégradation aussi rigoureux que les tests nominaux \\
12 & La carte reste, mais montre la jointure &
  Empreinte colorée par le test, panneau de vulnérabilité \\
13 & Le géocodage est un composant produit &
  Tests dédiés : arrondissements, seuil de confiance, saturation \\
\bottomrule
\end{tabularx}
\end{table}
